\documentclass{ltjsbook}
\usepackage{docmute} 
\usepackage{../../myStyle} 

\begin{document}
\VerbatimFootnotes
% \chapter{いんたーみっしょん；Stanの概略と環境の準備について}
\chapter{Introduction; About the outline of Stan and preparation of the environment}
\label{PrepareStan}

% このセクションでは，確率的プログラミング言語Stanを導入するにあたっての，周辺知識を解説します。授業中に解説するものではありませんし，ここに書かれていることのすべてを理解していないとStanを使えないというわけではありませんが，導入や利用にはトラブルやエラーも多く，その際に前提知識，周辺知識があるとないとでは理解度が大きく異なります。「わからなくても，なんとなく動いた」という状態で受講し続けるのは自身のためにならないだけでなく，\textbf{面白くない}です。知識があってはじめて価値がわかるということもありますので，共用としてご一読いただければと思います。
In this section, we will explain the background knowledge needed to introduce the probabilistic programming language Stan. It is not something that will be explained during the class, and you do not need to understand everything written here to use Stan. However, there are many troubles and errors in the introduction and use, and the level of understanding differs greatly whether you have prior knowledge and background knowledge or not. It is not useful, let alone interesting, to continue attending in a state of "I don't know, but it somehow worked". There are also times when you understand its value only when you have knowledge, so I hope you will read it for your own sake.


% \section{はじめに}
\section{Introduction}
% この授業では統計言語としてのR，それを使う統合開発環境としてのRStudio，そして\keyterm{確率的プログラミング言語}{かくりつてきぷろぐらみんぐげんご}{stocastic programming language}としてのStanを利用します。大学でのPCルームの利用にはこれらの環境がすでにある程度準備されていますが，最新バージョンではありませんので，自分のPCに環境を準備することを強く推奨します。今後の研究室配属やその後の卒業研究などでも活用することになると思いますので，まずは自身のPC環境に準備することを考えてみてください。
In this course, we will use R as a statistical language, RStudio as an integrated development environment, and Stan as a \keyterm{stochastic programming language}. Although these environments are partially prepared in the PC rooms at the university, they are not the latest versions, so we strongly recommend setting up the environment on your own PC. We believe you will utilize them in future lab assignments and graduation research, so first consider preparing them on your own PC environment.

% この付録資料は，これらの環境を準備する方法について解説するものですが，提供されるソフトウェアのバージョンや対応する環境などは日々発展するものですので，必ずしも最適な情報提供になり得ません。基本的な情報は提供いたしますが，執筆時\footnote{この記事は2022年11月10日に加筆修正されました。}での情報であることも多く，より新しい情報についてはインターネットなどでキーワード検索を行なって，より良いものを探してみてください。
This appendix explains how to prepare these environments, but the versions of the software provided and the corresponding environments evolve daily, so it may not always provide the optimal information. We will provide basic information, but a lot of the information is as of the time of writing\footnote{This article was revised on November 10, 2022.} Therefore, for newer information, please conduct keyword searches on the Internet and find better ones.

% PCの操作が苦手に思っている人のために付言しますが，うまくいかなかったからといって癇癪を起こしたり，諦めたりしてはいけません。相手は機械ですから，命じられたことを愚直に実行しているだけです。また文房具のようなものですから，恐れて嫌うのではなく，愛して可愛がってあげることが肝要です。もしみなさんの中に，お持ちのPCに名前をつけていない人がいるようでしたら，今すぐ命名することをオススメします。パソコンが，と思うと腹が立ちますが，わたしのXXXちゃんが，と思うとエラーも「ちょっとご機嫌斜めなのかしら」と受け入れやすくなります\footnote{ちなみに私が初めて手に入れたノートPCには「秀吉」という名前をつけました。Macに変わってからは代々数学者の名前をつけることにしています。私はいま，GaussちゃんやHermiteちゃんを持ち歩いています。}。
To those who find operating a PC difficult, I would like to say this; don't get upset or give up if things don't go well. Remember, your partner is a machine, and it merely follows the given instructions obediently. Since it functions like a piece of stationery, it is crucial not only to fear or dislike it, but also to love and cherish it. If there is anyone among you who has not yet named your PC, I highly recommend doing so right now. When your PC makes a mistake, it's easy to get angry if you think of it as just a computer, but if you think of it as 'my little XXX', even errors can more easily be seen as 'maybe it’s just in a bad mood'. As a footnote, the first notebook PC I ever owned was named 'Hideyoshi', and since switching to Mac, I have been naming them after mathematicians. Currently, I'm naming my devices as 'Gauss-chan' and 'Hermite-chan'.

% また，困った時は教員，TA，先輩などに相談することが重要ですが，その際には症状や経緯を正確に伝えることが必要です。ペットを病院に連れていくときに「何もしてないのに勝手に病気になった。治してほしい」といわれても獣医さんは困ると思います。\textbf{普段の様子や症状についての丁寧な報告}を心がけてください。とくにPCは機械ですので，何もしていないのに壊れるということはありません。自分が自覚していないことであっても，「なにかをした」から様子がおかしくなるのです。自分はそんな大それたことはしていない\footnote{たとえば計算途中だけど時間が来たので電源をオフにした，といったことでも，内部ファイルにアクセしているときの中断であれば，十分にPCを破壊することができます。}，と思っても必ず何かトリガーがあったはずなのです。どこをクリックしたか，どういうコマンドを入れたかという\textbf{履歴をしっかりと把握し，丁寧に報告する}ことを心がけてください。
Also, while it's important to consult with teachers, teaching assistants, seniors, etc. when you're in trouble, it's necessary to accurately convey the symptoms and circumstances. Imagine if you take your pet to the vet and say, "It got sick on its own without doing anything. I want you to cure it." It would be tough for the vet, wouldn't it? Please make sure to give a detailed report on your usual conditions and symptoms. Especially with PCs, as they are machines, they don't break without doing something. Even if you are not aware of it, they behave strangely because "something has been done". You might think you haven't done anything bad, but there must have been a trigger for it. Please make a habit of understanding what you clicked on, what kind of commands you put in, the history of your actions, and report them carefully.

% \section{Stanの位置付け}
\section{Positioning of Stan}

% \keyterm{確率的プログラミング言語}{かくりつてきぷろぐらみんぐげんご}{stocastic programming language}とは，確率モデルを記述し，データと合わせることで統計的推論をしてくれるコンピュータ言語のことです。Stanができる前は，\keytermL{BUGS}{BUGS}や\keyterm{JAGS}{JAGS}というものがありましたが，BUGsは開発が終わってしまいました\footnote{JAGSはまだ開発が続いているようで，Rから使うこともできます。}。今，こうした言語はStanが最先端だといって間違い無いでしょう。これらの言語は，具体的には確率モデルとデータの関係を記述することで，事後分布からの乱数を生成するというものです。以下ではこの言語の利用形式について解説します。
A stochastic programming language is a computer language that describes a probabilistic model and provides statistical inference by combining it with data. Before Stan was developed, there were languages like BUGS and JAGS, but BUGS' development has ended. JAGS still seems to be under development and can be used from R. Currently, it is safe to say that Stan is at the cutting edge of such languages. These languages specifically describe the relationship between a probabilistic model and data to generate random numbers from the posterior distribution. The following will explain the usage pattern of this language.

% \subsection{コンパイラとインタプリタ}
\subsection{Compiler and Interpreter}

% プログラミングはコンピュータを動かすための仕様書を書くことです。ここで「コンピュータ」というのは計算機という意味だと思ってください。もちろんコンピュータは数字の計算だけでなく，音楽を奏でたりゲームを楽しませてくれたり，と色々なことができるのですが，その背後にあるのはとにかく0/1の数値演算です。0か1かという数字がどうして映像や音声，通信になるのか，と疑ってしまうほど異なっているようですが，それでもすべて数字の計算から構成されています。
Programming is about writing specifications for moving a computer. By "computer", think of it as a calculator. Although a computer can do various things like play music or entertain games, behind all these is the operation of 0/1 numbers. How the numbers 0 or 1 turn into images, sound, and communication may seem questionable, yet everything is still constructed from number calculations.

% コンピュータができ始めたごく初期の頃は，これを動かすのに0と1からなる数字の羅列で仕様書を用意していました。流石にそれではわかりにくく表現力にも乏しいので，次にできたのがマシン語とよばれる16進数による表現でした。この段階でも，知らない人にとっては暗号か，意味のある連なりに思えない文字列にすぎません。画面に線を引いて欲しい時に\texttt{line}という命令で伝えられるようになって，やっと普通の人間にも意味がわかるレベルになってきます。このように，人間がわかるレベルでコンピュータに命令が伝えられるような言語のことを\keytermL{高級言語}{こうきゅうげんご}と言います。高級言語は人間寄りで，直接コンピュータが理解できませんので，この高級言語を機械の言葉に翻訳するためのアプリケーションを介在させます。それがいわゆる\keyterm{プログラミング言語}{ぷろぐらみんぐげんご}{programming language}と呼ばれるものです。
In the very early days when computers were first developed, the specifications for operating them were made up of a string of numbers consisting of 0s and 1s. Needless to say, this method was difficult to understand and lacked expressive capabilities, so the next thing that was developed was an expression using hexadecimal, also known as machine language. Even at this stage, for those who are unfamiliar, the string of characters would seem like nothing more than a code or meaningless sequence. When you want a line drawn on the screen and this can be conveyed with a command such as \texttt{line}, then it finally reaches a level where even ordinary people can understand its meaning. Such languages that allow instructions to be conveyed to computers at a level comprehensible to humans are called \keytermL{high-level languages}. High-level languages are more human-friendly, but can't be directly understood by computers, so an application is used to translate these high-level languages into machine language. This is what is commonly referred to as a \keyterm{programming language}.

% プログラミング言語は，古くはBASIC，PASCAL，FORTANなどと呼ばれる書式のものがあり，その後出てきたC言語，その改良版であるC++言語\footnote{この\verb|++|という書き方は，C言語特有の「変数に1加える」という表記法であり，C言語の改良版という意味でC++と命名されています。読み方はシープラスプラスですが，シィプラプラの愛称でも知られています。}などが有名です。他にもJAVAやObject C，Pythonなどが有名ですね。Rも言語の一種で，統計解析に特化したものです。また，この授業で扱うStanという\keyterm{確率的プログラミング言語}{かくりつてきぷろぐらみんぐげんご}{stocastic programming language}は，確率モデルの分析に特化したものです。みなさんがデータサイエンス業界に進むのであれば，PythonやRを使えるようになっておくといいでしょう。これらの言語の特徴の1つは，パッケージを追加することで機能が拡張できることにあります。とくに\keytermL{機械学習}{きかいがくしゅう}系(AIなど)のパッケージが豊富なのがPythonです。またこれらの言語は基本的にフリーソフトウェアであり，導入にお金がかからないところもポイントですね。タダで始められるおもちゃみたいなものです！\footnote{余談ですが，Scratchや任天堂Switchの「ナビつき! つくってわかる はじめてゲームプログラミング」などにみられる，GUIベースのプログラミング言語もあります。これはプログラミングの要素がアイコンやキャラクターになっていて，それらを組み合わせることで計算させるという方法をとっています。その操作のしやすさ(ミススペルがない！)から，小学生などにも導入できる素晴らしい取り組みです。同様の発想は，LOGO言語などにもみられました。}\footnote{PerlやJava Scriptなど，Webブラウザ上で動くプログラミング言語もありますが，これはブラウザ上での操作に限定されており，数値計算にはあまり向かないのでここには取り上げていません。さきほどあげたJAVAとJAVA Scriptは別物なので注意してください。}\footnote{GUIアプリケーションを作ることに特化した言語もあります。数値計算と違って，マウスがどこでクリックされたか，と言った「動き」を契機としてプログラムが走る，イベントドリブン型の言語で，Microsoft社のVsial Basicなどが有名です。これはMicrosoft Office製品の中に埋め込むこともでき(VBA)，これを使うとExcelでも高度な統計解析ができたりします。Excelを使った統計ソフトウェアとして代表的なものにHADがあります\parencite{HAD}。他にもDelphiなどの言語がありました。}\footnote{\parencite{Language2016}には，数え方にも癖がありますが，117ものプログラミング言語が紹介されています。}
Programming languages have been around for some time, including formats such as BASIC, PASCAL, and FORTAN. This was followed by the C language and its improved version the C++ language. The "++" in C++ is a notation unique to the C language meaning "to add 1 to a variable", thus showing that C++ is an improved version of C. It's pronounced "see-plus-plus" and is also known by its nickname, "see-pra-pra". Other famous ones include JAVA, Object C, and Python. R is another type of language, specialised for statistical analysis. Additionally, a stocastic programming language called Stan, which will be discussed in this course, is specialised for analyzing probabilistic models. If you plan on advancing into the data science industry, it would be beneficial to become proficient in Python or R. One feature of these languages is that functionality can be expanded through additional packages, particularly in Python, which has an abundance of machine learning packages including AI. Moreover, these languages are generally free software, so the fact that they don't require any monetary investment is also a point to note. It's like being able to start with a free toy!

As an aside, there are also GUI-based programming languages like Scratch and Nintendo Switch's "Game Builder Garage". In these languages, programming elements become icons and characters, and calculations can be made by combining them. The ease of operation (no misspelling!) of these languages is a great initiative that can also be introduced to elementary school students. Similar concepts could also be seen in the LOGO language.

There are also programming languages that run on web browsers, such as Perl and JavaScript, but these are limited to browser-based operations and are not suitable for numerical calculations, so they are not included here. Please note that JAVA and JavaScript are not the same thing.

There are also languages specialised for creating GUI applications. Unlike numerical calculations, these event-driven languages, such as Microsoft's Visual Basic, can run programs based on "movements" such as where the mouse was clicked. It can also be embedded in Microsoft Office products (like VBA), and using it can allow for advanced statistical analysis in Excel as well. HAD is a leading statistical software that uses Excel. There are also other languages such as Delphi.

According to a source, there are 117 programming languages, although the counting method can be idiosyncratic.

% プログラミング言語の分類には，その設計思想や書き方などを基準にすることもできますが，実行方法を基準にするものとして\keytermL{インタプリタ}{いんたぷりた}型と\keytermL{コンパイラ}{こんぱいら}型とに分けることができます。
The categorization of programming languages can be based on their design concepts and the way they are written. However, they can also be divided into two types based on their execution method: \keytermL{interpreter}{interpreter} type and \keytermL{compiler}{compiler} type.
% インタプリタ型は，interpret，つまり翻訳型です。これは毎回の命令を機械語に翻訳して実行していきます。Rはインタプリタ型言語で，毎回の命令を逐一翻訳して実行していきます。Rを実行するとき，コンソールに\verb|>|というマークが出ていますよね。これはRが聞き耳を立てている状態，入力待ちの状態なのです。ここに命令を書くと(たとえば\verb|2+3|と書く)，Rはすぐさま答えを返してきます(たとえば\verb|[1] 5|と返す)。基本はこのように一問一答，ひとつひとつの命令を毎回機械語に翻訳して実行し，その答えを出力するという手続きをとっています。もちろん私たちの分析プログラムは，一行で書ける簡単なものではありませんから，複数行に渡る長々としたファイルを書くことが多いでしょう。こうしたファイルはプログラムともいわれますし，一行一行の命令のことを\keyterm{スクリプト}{すくりぷと}{script}ということもあります。プログラムファイルあるいはスクリプトファイルを開いて，一行ずつ実行するのがRの基本的なスタイルであり，RStudioはスクリプトファイルを編集する\keyterm{エディタ}{えでぃた}{editor}がセットになっているのが便利な点なのでした。
The interpreter type is interpretive, i.e., it translates each instruction into machine language and executes it. R is an interpreted language and translates and executes each instruction one by one. When you run R, you see a '>' mark on the console, don't you? This is R in it's active state and waiting for input. If you write an instruction here (e.g., writing '2+3'), R will immediately return an answer (e.g., returning '[1] 5'). The basic procedure is to ask and answer one by one, translating each instruction into machine language and executing it, and outputting the answer. Of course, our analysis programs are not simple enough to be written in one line, so we often write long files across multiple lines. Such files can be referred to as 'programs', and each line of instruction can also be referred to as a 'script'. The basic style of R is to open a program file or script file and execute it line by line, and the convenience of RStudio is that it comes with an editor that allows you to edit script files.


% これに対して，インタプリタ型は，スクリプトファイルをまとめて機械語に翻訳(コンパイル)し，実行ファイルというのを別途作ります。その上で，実行ファイルを実行すると計算が進められるのです。どうしてこんな手間がかかることをするのでしょうか。ひとつには，ひとつひとつの命令分を逐一翻訳しているのでは，実行スピードが遅くなるということが挙げられます。命令を逐一聞き耳を立てて待ち，全ての命令を聞き終わるまで途中の指示を記憶しておいて，命令文が終わってはじめて行動に移す，というのでは記憶容量も時間もかかるわけですね。これに対して，一連の計算命令文書を一括で渡すことができれば，機械は内部的に効率よく計算手順を配置して実行できます。コンパイラ型は計算スピードが速いのです。ちなみに実行ファイルは機械がわかる言葉に書き換わっているので，人間が見ても意味がわかりません。また，OSやCPUに理解できる形に書き換えていますので，MacOSの実行ファイルをWindowsで実行する，ということはできません。翻訳後の言語が違うからです。スクリプトファイルは環境を超えて共有できますが，それを翻訳したり最適化したりする仕組みは，個別の環境に依存します\footnote{ちなみにRやRstudioはもちろん，WordやExcelなどのオフィスアプリ，果てはOSそのものも，コンパイルされた実行ファイルをPCが計算して実行しているに過ぎません。}。
In contrast, interpreted programming languages compile the entire script file into machine language and create a separate executable file. The calculations are carried out by running this executable file. Why would we go through such a process? One reason is that when each instruction is translated individually, the execution speed becomes slow. If you listen to each instruction one at a time, wait for the machine to catch all instructions and keep the interim instructions in mind until all instructions are listened to, and then finally start moving according to these instructions, it will take time and memory capacity. In contrast, if we can deliver a series of computational instructions all at once, the machine can efficiently organize the computational procedures internally and execute them faster. This is why compiler languages execute faster. By the way, the executable file is translated into a language that the machine can understand, so it is meaningless when a human tries to read it. It's also written in a form that OS and CPU can understand, so you can't run a MacOS executable file on Windows because the language is different after translation. Script files can be shared across different environments, but the system for translating or optimizing them depends on the individual environment. By the way, R or Rstudio, and even office applications like Word or Excel, and even the OS itself, are merely executing compiled files on the PC.

% Stanはコンパイラ型です。Stanのファイルを読み込んで，機械にわかるように「事後乱数生成命令文」に翻訳して実行します。確率モデルから乱数を生成するのは非常に高度な計算機能ですから，コンパイルしてまとめておくことでスピードアップの効率が断然良くなるわけです。コンパイラ型の利点はこのスピードにありますが，欠点は「命令に変更やミスがあれば翻訳し直さなければならない」というところです。あり得ない命令を出しているようであれば，コンパイルの段階で「おかしな文法だよ」とエラーを出して翻訳を止めてくれます。翻訳が通れば良いかというとそうではなく，翻訳はあくまでも文法上の手続きであって，数値や内容に間違いがあっても翻訳自体は完了させることができるのです。翻訳(コンパイル)には少し時間がかかりますから，文法上も内容上も間違いのない命令分をしっかり準備しておくことが重要になってきます。
Stan is a compiler type. It reads Stan files, translates them into "posterior random number generation commands" that machines can understand, and executes them. Generating random numbers from probability models is a highly advanced computational function, so by compiling and bundling them together, the efficiency of speed-up becomes drastically better. The advantage of a compiler type lies in its speed, but the disadvantage is that "if there are changes or mistakes in the instructions, they must be translated again". If it seems that you are issuing impossible commands, it will stop the translation by issuing an error that says "the syntax is strange" at the time of compilation. Just because the translation can get through doesn't mean it's okay. Translation is just a grammatical procedure, and even if there are mistakes in numbers or content, the translation itself can be completed. Translation (compiling) takes some time, so it's important to prepare instructions that are free of grammatical and content errors.
% \subsection{ファイルのやり取りについて}
\subsection{About File Exchange}
% ここで，R/RStudioをつかってStanを使う際の，ファイルのやり取りを見ておきましょう(図\ref{fig::pp_01})。
Let's take a look at the file exchange when using Stan with R/RStudio here (Figure \ref{fig::pp_01}).

% StanのコードはRのコードの中に書くこともできますが，別のファイルに保存しておくのがベターです。Stanファイルは拡張子を\verb|.stan|とすることが一般的です。このファイルの中身はプレーンテキストでコードを書いていますから，一般的なエディタ\footnote{エディタとはASCIIファイル，いわゆるプレーンテキストを編集するアプリケーションの総称で，OSにデフォルトで「メモ帳」などの名称で含まれています。デフォルトのアプリは文字を読み書きできるだけですので，もう少し発展的，あるいは便利な拡張機能を持っている専用のアプリを使うことが一般的です。筆者は最近もっぱらVS Codeというエディタを利用しています。他にもMacOSでしたら，「mi」というアプリを好んで使っていました。Windowsでしたら「秀丸」というアプリが有名です。Ubuntuにはgeditというアプリがデフォルトで入っていますし，歴史やユーザの多さではvimやemacsなどがあります。Linux界隈で，vim派かemacs派かの話題は，きのこたけのこ戦争ばりに激しい戦いになります。}で編集できます。RStudioのエディタ画面も必要十分なエディタ機能を持っていますので，RStudioをエディタとして利用するのがいいでしょう。
You can write Stan code in R code, but it's better to save it in a separate file. Stan files are typically given the extension \verb|.stan|. The contents of this file are written in plain text code, so any general editor can be used to edit it. An editor is a generic term for an application that edits ASCII files, or plain text, and is included by default in the OS as "Notepad" or something similar. The default application can only read and write characters, so it's common to use a dedicated application that has a bit more advanced or convenient functions.  Recently, I have been mainly using an editor called VS Code. I also used to like using an app called "mi" on MacOS. For Windows, there's a popular app called "Hidemaru". Ubuntu comes with a default app called gedit, and there are also others like vim and emacs, which have a long history and large user base. In the Linux community, the topic of whether you belong to the vim camp or the emacs camp can lead to intense debates comparable to the mushroom and bamboo shoot war. The editor screen of RStudio also has a sufficient editing function, so it would be good to use RStudio as an editor.


% RStudioでFile > New Fileとして新しいファイルを開くときに，Stanファイルとして開くことが可能です。Stanファイルとして開くと，初心者向けの配慮からか，最初からちょっとしたコードがすでに書かれています。コードの書き方を示すサンプルコードなのですが，実際にこのコードを使って何かすることはありませんので，中身は全部削除してしまいましょう。それでもRStudioはその画面がStanファイルのものだということは認識していますから，そこで書いたファイルはStanファイルとして扱ってくれます。ここでまちがって新しくR Scriptで開いてしまった，というときに，保存するときだけ\verb|.stan|をつけておけばいいか，という対応をすると失敗します。RStudioでは拡張子をあえて表示していませんので，ファイル名として\verb|hoge.stan|と書くと実際には\verb|hoge.stan.R|というファイル名になっています。拡張子は最後のピリオドで判断されますので，これはRファイルなのです。このミスを防ぐために，ファイル名にピリオドは使わないこと，ファイルの種類を変える時はエディタ画面の右下でファイル種別を選択してください(図\ref{fig::pp_02})。
When opening a new file through File > New File in RStudio, it is possible to open it as a Stan file. When you open it as a Stan file, maybe consideration for beginners, a little bit of code is already written at first. This is sample code to show how to write the code, but you won't actually use this code to do anything, so let's delete all of its contents. Even then, RStudio recognizes that this screen belongs to a Stan file, so the file you wrote there will be treated as a Stan file. If you accidentally open a new R Script here and save it with \verb|.stan| appended, that response will fail. In RStudio, because the extension is deliberately not displayed, if you write \verb|hoge.stan| as a file name, the actual file name becomes \verb|hoge.stan.R|. Since the extension is determined by the last period, this is an R file. To prevent this mistake, do not use periods in file names and select the file type from the bottom right of the editor screen when changing the file type (Figure \ref{fig::pp_02}).

\begin{figure}[ht]
	\centering
	\includegraphics[keepaspectratio,width=0.8\linewidth]{../images/prepare/fileStyle.png}
	\caption{ファイル種別の変更}
	\label{fig::pp_02}
\end{figure}


% RStudioを使ってStanを書くと便利なことがいくつかあります。ひとつは強調表示機能で，ブロックや関数名などStanで使うことが決まっている専門用語は色やフォントが変わって表示されます。\verb|transformed parameters|のように長い専門用語を書くときはミススペルが心配ですが，書いた後で強調されなければスペルミスがある，ということがわかります。もうひとつの利点は文法チェックの機能です。コンパイル型なので，Stanに書かれた命令文は実行前に一括で翻訳されますが，スペルミスや型の違いなど，文法的に間違っているところがあればこれも自動的にチェックして警告記号がでます。またStanファイルが開かれているペインの左上に\verb|Check|というボタンがあり，これを押すことでファイル全体の文法チェックをしてくれます。もし文法上の問題がなければ，「hogehoge.stan is syntactically correct.(hogehoge.stanというファイルは文法的には正しいですよ)」というメッセージがコンソールに表示されます。逆にいうと，これが表示されないというときは何か間違っているので，コンパイルに進む前に修正しましょう。
There are several advantages to using RStudio to write Stan. One is the highlighting feature, which changes the color and font of specialized terms used in Stan, such as blocks and function names. When you are writing a long professional term like \verb|transformed parameters|, you may be worried about misspelling, but you will know there is a spelling mistake if it does not get highlighted afterwards. Another advantage is the grammar checking function. Stan's commands are translated in bulk before they are executed because it's a compiled language. If there are grammatical errors, such as spelling mistakes or type differences, these are automatically checked and a warning sign is displayed. Also, there is a \verb|Check| button at the top left of the pane where the Stan file is opened, and pressing this will check the entire file for grammar. If there are no grammatical problems, the console will display a message stating "hogehoge.stan is syntactically correct." Conversely, if this message is not displayed, something is wrong, so make corrections before proceeding to compile.

% さてこのようにしてStanファイルを準備します。また分析用のデータセットが外部ファイルにある場合などは，これまで同様，当該プロジェクトフォルダ内に置いておくと良いでしょう。これらはいずれも，Rのコードで呼び出して使うことになります(図\ref{fig::pp_01}の1.と2.のステップ)。Rでは，Stanファイルやデータファイルを読み込んで，これをPC内部にあるStanに渡します(図\ref{fig::pp_01}の3.のステップ)。StanはRと違う言語であり，Stan独自の計算機能で計算してくれるわけですから，RとStanの橋渡しが必要になります。Rではこれを\verb|cmdstanr| や\verb|rstan|というパッケージを経由して行います。これらのパッケージがStanを呼び出してくれるわけです\footnote{このように，Stanは独立した計算環境であり，R以外のアプリケーションから呼び出して使うことが可能です。Pythonから使いたい場合は，PyStanというパッケージ経由で，Juliaから使いたい場合はStanJl，コマンドラインから使いたい場合はcmdstanといったように，いろいろなルートがあります。}。
So we prepare the Stan file like this. Also, if the dataset for analysis is in an external file, it would be a good idea to place it in the project folder as before. Both of these will be invoked and used in R code (steps 1 and 2 in Figure \ref{fig::pp_01}). In R, we load the Stan file and data file and pass them to Stan inside the PC (step 3 in Figure \ref{fig::pp_01}). Since Stan is a different language and computes with its own calculation function, we need a bridge between R and Stan. In R, this is done through packages called \verb|cmdstanr| and \verb|rstan|. These packages invoke Stan\footnote{As such, Stan is an independent computing environment and can be invoked and used from applications other than R. If you want to use it from Python, you can do so through a package called PyStan, and if you want to use it from Julia, you can do so through a package called StanJl. If you want to use it from the command line, there are various routes such as cmdstan.}.

% データと確率モデルの関係を記したStanファイルと，データのそのものを受け取ったStanは，コンパイルしてPC内に乱数発生装置を作ります(図\ref{fig::pp_01}の4.のステップ)。事後分布の関数を書き出すのではなく，その形状とそれを代表する乱数を作る状態をもつわけです。Rの指示をうけて，そこから必要な数だけ乱数を作り出します(図\ref{fig::pp_01}の5.のステップ) 。Rはこれを受け取って，統計解析を始めるわけです。Rは統計に特化した言語ですから，集計や可視化などの作業はお手のもの。事後分布から得られた大量の乱数は，事後分布の特徴を反映した大量のデータセット，事後分布の具体例たちになっているのでした。
Stan, which receives both a Stan file that lists the relationship between data and probability models, and the data itself, will compile and create a random number generator in the PC (the 4th step in Figure \ref{fig::pp_01}). It has the state of creating the shape of the posterior distribution and the random numbers that represent it, rather than writing out the function of the posterior distribution. Upon the direction of R, it generates a necessary number of random numbers from there (the 5th step in Figure \ref{fig::pp_01}). R, accepting this, starts statistical analysis. Since R is a language specialized in statistics, tasks such as aggregation and visualization are its strengths. A large amount of random numbers obtained from the posterior distribution have become a large data set reflecting the characteristics of the posterior distribution, concrete examples of the posterior distribution.

% 分析の進め方は以上のようになります。Rでデータを整形し，Stanのファイルも別途用意して，RからStanを呼び出す，Stanか返してきたデータをRで解析する，ということです。Rの拡張子は\verb|.R|で，Stanの拡張子は\verb|.stan|，データファイルの拡張子は\verb|.csv|などでしょうか。いずれにせよ，複数の種類のファイルをやり取りしながら進めるので，相互の関係についてしっかり理解しておく必要があります。
The procedure for the analysis is as follows: you shape the data in R, prepare separate Stan files, call Stan from R, then analyze the data that Stan returns in R. The extension for R is \verb|.R|, for Stan, it's \verb|.stan|, and for the data files, it might be \verb|.csv|. Regardless, you need to understand the relationship between the multiple types of files you'll be exchanging, as you will proceed by exchanging them back and forth.


\begin{figure}[ht]
	\centering
	\includegraphics[keepaspectratio,width=0.8\linewidth]{../images/prepare/R_stan.png}
	\caption{R/RStudioとStanファイルのやり取り}
	\label{fig::pp_01}
\end{figure}

% \subsection{2つのルート}
\subsection{Two Paths}
% ここではRからStanを呼び出す2つのルート，具体的には\verb|rstan|パッケージと\verb|cmdstanr|パッケージについて説明します。
Here, we will explain two ways to call Stan from R, specifically, the rstan package and the cmdstanr package.
% この2つのパッケージは，いずれもStanを呼び出してつかう，RとStanの間を取り持つインターフェイスの役割を持ったパッケージです。間を取り持つだけなので，Stanファイル自体は同じであっても構いません。同じStanファイルを\verb|rstan|から呼んでも，\verb|cmdstanr|から呼んでも，結果は同じです。ただしインターフェイスが違うので，結果の扱い方が変わってきます。他にも違うところがいくつかありますので順に説明していきますが，結論からいうとこれから\footnote{執筆時点，2022-11-14現在}は\verb|cmdstanr|のほうを使うほうがお勧めです。
These two packages, both of which call and use Stan, serve as interfaces between R and Stan. Since they only mediate, it doesn't matter if the Stan file itself is the same. Whether you call the same Stan file from \verb|rstan| or \verb|cmdstanr|, the result is the same. However, the interface is different, so how the results are handled changes. There are also a few other differences, which I will explain in order, but the conclusion is that I recommend using \verb|cmdstanr| from now on (as of the time of writing, November 14, 2022).

% 歴史的には\verb|rstan|のほうが古くから存在します。このパッケージも最初の頃は導入に一苦労するものでしたが，2016年ごろからCRANを通じて配布されるようになりました。要するに，他のパッケージと同じように，\verb|install.packages("rstan")|とRで書くだけで導入できるようになったのです。このようにして導入するとわかりますが，\verb|rstan|は多くの依存パッケージがあります。StanはコンパイルするのにOSに応じたC++コンパイラを借用しますが，RからC++を呼び出すパッケージやRとC++の間に入るStanHeaderと呼ばれる緩衝材など，関連する多くの環境を整えてやっとRからStanへのルートが通じるのです。ユーザにとってはそのような苦労は知ったことではない，という感じですが，「間に多くの調整が入る」ということがエラーの温床になっていました。つまり，いくつかの内部的ステップのどこかでバグがある，互換性がない，というようなことがあると，Stanファイルが同じでも，バージョンが変わるとコンパイルできないということがよくあったのです。Stan本体の方も開発が盛んですから，Stanの新しいバージョン，新しい機能を使いたいと思っても，パッケージが対応するのを待たなければなりません。パッケージが対応しても，途中の媒介パッケージが追いついていなければバグになったりします。これらを使うRやOSもアップデートしていきますから，その度に「動かなくなる」ということはよくありました。研究の再現性が問題になる昨今ですが，自分の研究とデータであっても，自分の環境で再現できないということが少なくなかったわけです\footnote{実際RやOSのアップデートに関わるエラーは大変で，筆者は一度OSをアップデートした後半年ほどStanが使える環境が作れず，研究が頓挫したことがあります。そのために別途，PCを購入し直したぐらいです。}。そうなると，今動いているからそれでいい，と環境のアップデートを控えるようになりますが，ご存知の通りOSやアプリのアップデートはバグやセキュリティホールの修正など，安全につかうための基本的な機能に関わってきますから放置もできない，という問題があるわけです\footnote{OSのアップデートなど，通知が来ても「今支えているからいいや」と無視する人もいますが，できればOSやアプリは最新版であるほうが良いのです。セキュリティホールの問題などを放置しておくと，インターネットを介して自分のPCの中身が全部見られて公開される危険性があります。Stanしか使わないネットに繋がないPCというのがあればいいのかもしれませんが，文房具としてのPCはもっと一般的な用途がありますから，そうもいってられませんよね。また，一昔前のOSアップデートは，アップデートの失敗でPCが動かなくなるというようなこともありましたから，アップデートが公開されてすぐに対応するのは「人柱」などと揶揄されていましたが，最近はそういったこともほとんどありませんので，安心して常に最新の状態にしておいてください。}。
Historically, \verb|rstan| has existed from an early time. Initially, it was a struggle to install this package, but it became available through CRAN around 2016. In other words, like other packages, it became available by simply writing \verb|install.packages("rstan")| in R. As you can see when you install it, \verb|rstan| has many dependent packages. While Stan borrows a C++ compiler according to the OS to compile, it takes many related environments, such as packages to call C++ from R and a buffering material called StanHeader that sits between R and C++, to finally establish a route from R to Stan. Users might not care about these struggles, but the fact that "many adjustments are made in the process" has often led to errors. What I mean is, even if the Stan file is the same, changes in version can often result in inability to compile due to bugs or incompatibility at some of the internal steps. Stan itself is also actively developed, so even if you want to use the new version or features of Stan, you have to wait for the package to support it. Even if the package does support it, there could be bugs if the intermediary packages don't keep up. Since both R and the OS that use these are also updated, it often happens that "it stops working" each time. There is a growing problem of reproducibility in research, but even with my own research and data, it was not uncommon to be unable to reproduce it in my own environment\footnote{Actually, errors related to updates of R or the OS are serious, and I had a case where I couldn't create an environment where I could use Stan for about half a year after updating the OS, which stalled my research. So much so that I bought a new PC for it.}. That makes you hesitant to update your environment because it's working now, but as you may know, updates to the OS and apps are involved in basic functions for safe use, such as bug fixes and security hole patches, so you can't just leave it\footnote{Some people ignore notifications about OS updates, thinking "it's working now, so it's fine", but it's better to have the latest version of the OS and apps if possible. If you leave security holes unattended, there is a risk that the contents of your PC may be exposed and published through the Internet. While it might be nice to have a PC that doesn't connect to the net and only uses Stan, PCs as stationery have more general uses, so you can't just say that. Also, in the old days, OS updates could sometimes cause the PC to stop working if the update failed, and people who responded immediately after an update was released were sometimes ridiculed as "human pillars", but recently there are hardly any such problems, so please keep your system up to date with confidence.}.


% このような不安定な環境になる原因は，RとStanの間にあるさまざまな障壁と，それをなんとかして調整しようという媒介パッケージの存在でした。そこでなるべくシンプルにRとStanをやりとりする方法はないものか，ということで考えられたのが\verb|cmdstanr|です。これはコマンドライン\footnote{MacOSやLinuxでは端末，ターミナル，Windowsではコマンドプロンプトなどと呼ばれる，PCに直接コマンドで命令を出すプリミティブな実行環境です。}からStanを呼び出す\verb|cmdstan|をRから呼び出して使おう，というものです。StanをRになんとか取り込んで動かす，というのではなく，Stanの部分はOSのプリミティブな環境にお任せ，Rは結果をもらうだけ，というシンプルな設計にしようというわけです。たとえば\verb|rstan|はStanの計算結果をRオブジェクトとして取り込みますが，\verb|cmdstanr|は実はMCMCの結果はcsvファイルのような形で吐き出しており，それを読み込んで使います。御用聞きや仲介業者を介して情報をもらうのではなく，現地で直接買い付けしてくるようなイメージでしょうか。このように間に挟むものを減らすことによって，不安定な要素をなくしたわけです。
The cause of such an unstable environment was the various barriers between R and Stan, and the existence of mediator packages that somehow coordinate them. This prompted the consideration for a simpler way to interact between R and Stan, hence the inception of \verb|cmdstanr|. This is about using \verb|cmdstan|, which calls Stan from the command line\footnote{This is a primitive execution environment in which you give direct commands to your PC. It's known as a Terminal in MacOS and Linux, and Command Prompt in Windows.}, from R. Instead of trying to incorporate and run Stan in R, the idea was to leave the Stan part in the OS's primitive environment and have R just receive the results. This is a simple design idea. For example, while \verb|rstan| imports Stan's calculation results as R objects, \verb|cmdstanr| actually outputs the result of MCMC as a csv file or the like, and reads and uses it. It's kind of akin to directly purchasing information on-site rather than receiving it through a middleman or concierge. By reducing the intermediaries in this way, unstable elements are eliminated.

% \verb|cmdstanr|はStanとRの複雑な絡まりを排していますので，Stanの開発が進めばすぐにそれをRに届けることができます。たとえば2022/11/14現在で，公式にリリースされている\verb|rstan|のバージョンは2.21.1ですが，\verb|cmdstanr|が呼び出すstanのバージョンは2.30.1です。\verb|cmdstanr|のほうが新しいものに対応できていますね。このように，\verb|cmdstanr|がでてきてからはこちらの対応，反応が早く，また動作も安定的ですから，よりお勧めしやすくなっています。
The `cmdstanr` eliminates the complex tangle between Stan and R, so it can deliver the progress of Stan development to R immediately. For example, as of November 14, 2022, the version of `rstan` that is officially released is 2.21.1, but the version of Stan called by `cmdstanr` is 2.30.1. The `cmdstanr` is more compatible with newer things. In this way, since the advent of `cmdstanr`, its response and reaction have been faster and its operation has been more stable, which has made it easier for us to recommend it.

% これまでに出ているRでStanを使う方法を解説したテキストの多くは，\verb|rstan|をつかっていますので，今それを使うためには多少の読み替えが必要です。とはいえ，Stanのコード自体はそのまま使えるものがほとんどで，Rとのインターフェイス，やりとりの方法が違っているだけです。\verb|cmdstanr|には，結果を\verb|rstan|で得られたオブジェクトに変換する関数も含まれていますから，これらを使って変換すれば，以後のコードは\verb|rstan|準拠のものでもエラーになることはありません。このテキストでは第2版から\verb|cmdstanr|を中心にするように切り替えましたが，以前のコードでも本質的に問題はありません。
Many of the texts that have so far explained how to use Stan in R are using \verb|rstan|, so some reinterpretation is needed to use it now. However, most of the Stan code itself can be used as is, and only the interface and interaction methods with R are different. The \verb|cmdstanr| includes functions to convert the results into objects obtained with \verb|rstan|, so if you use these to convert them, you won't get errors even with \verb|rstan|-compliant code after that. This text has switched to focusing on \verb|cmdstanr| from the second edition, but there is no essential problem with the previous code.

% \section{導入の概略}
\section{Overview of Introduction}
% 環境の導入は，OSの種類によって変わります。WidowsなのかMacなのか，あるいはLinuxなのかといった違いに合わせて導入を考えてください。
The introduction of the environment varies depending on the type of OS. Whether it's Windows or Mac, or even Linux, please consider the introduction tailored to these differences.
% OSの種類で言えば，Linuxは全体的にCUI操作が主であり，堅牢かつ合理的な使い方ができるものです。如何せんマイノリティですので，ググってもあまり情報が出てこないところが玉に瑕です。
Speaking in terms of the type of OS, Linux is mainly operated by CUI and is something that can be used robustly and reasonably. However, since it is a minority, the only flaw is that not much information comes up when you search for it on Google.
% MacはLinuxベースのOSになっているので，比較的堅牢かつ合理的な使い方ができます。ハードウェアもOSも1つの会社(Apple)が作っているので，サポートもしやすいという利点がありますが，日本では少数派なのが残念なところです。
Mac is based on a Linux-based OS, so it can be used in a relatively robust and rational way. Both the hardware and the OS are made by one company (Apple), so it has the advantage of being easy to support, but it's a pity that it's a minority in Japan.
% Windowsは多数派のOSですが，OSメーカはソフトウェア会社でハードウェアが別会社なことが多く，利用される範囲は広いのですが，その分問題が生じたときにケースバイケースになりがちです。多くのハードウェアの違いを吸収するために最大公約数的な設計をするせいか，ソフトウェアデザインの統一性がないところが欠点です。ユーザ数は最も多いので，一生懸命調べたら同じような問題で困っている同じような機体の人に出会える可能性が高いのがせめてもの救いです。また光明としてWindowsOSもLinuxベースの仕組みを取り入れ始め(WLS)，比較的合理的な振る舞いができるようになったことが挙げられますが，OSのグレードによって導入のしやすさが異なります。
Windows is the majority OS, but the OS manufacturer is often a software company and the hardware is from a separate company. This makes it widely used, but it tends to be case-by-case when problems occur. The lack of uniformity in software design, possibly due to the design principle of accommodating a wide variety of hardware, is a disadvantage. Given that it has the most users, if you research hard you are likely to encounter someone with the same machine facing the same problem, which is a small consolation. On the bright side, Windows OS has begun to incorporate Linux-based mechanisms (WLS), enabling it to behave more logically, although the ease of implementation varies depending on the OS grade.

% ちなみにChromeOSやiOS, Androidなどタブレット・スマートフォンでの統計環境の利用は限界があります。これらは計算結果を利用するフロントエンドとしては非常に高機能なのですが，インタラクティブに使うことには不向きです。いわば書籍のように多くの情報を一方的に与えてはくれるのですが，ユーザが計算するといった働きかけの補助になるような(ノートとペンのような)使い方ができませんので注意してください。
By the way, there are limitations to using a statistical environment on tablets and smartphones such as ChromeOS, iOS, Android, etc. These are highly functional as front-end interfaces for using computational results, but they are not suitable for interactive use. In a way, they can provide a lot of information like a book, but they do not facilitate user calculations and involvement (as in the use of note and pen) so please be aware.

% さて，環境の導入として2つのルートを紹介しましょう。第一のものが最も基本的なアプローチで，自分の手元の環境を作り上げるというものです。これを\textbf{ローカル環境}での構築と呼びます。ローカル環境は自分だけの環境ということなので，自分の責任でもってしっかりと環境構築をできます。第二のアプローチとして，ある程度できあがった環境を丸ごと取り込む，\textbf{仮想環境}での構築という方法があります。PCのOSの中に別のマシン・別のOSをさらに憑依させて使うようなやりかたで，これができると取り込む環境は完成しているので細かい設定をする必要がありません。問題は「憑依させる」準備がいろいろ大変であるということです。また憑依させたPCはブラウザ経由でアクセスすることになることにも注意してください。
Now, let's introduce two routes to set up the environment. The first is the most basic approach, which is to create your own environment. This is called \textbf{local environment} construction. Since the local environment is your own space, you can carefully set it up under your responsibility. The second approach is to fully adopt a ready-made environment, which is called \textbf{virtual environment} construction. This is like inviting another machine or OS inside your PC's OS and using it. While importing, the environment is already complete, so there is no need to make detailed settings. The problem is that it's quite a hassle to prepare for the "possession". Also, be aware that you will access the possessed PC via a browser.
% 最後の方法は，\textbf{サーバの利用}です。これは環境を構築するのではなく，すでに誰かがどこかで構築してくれたサービスを受け取るという，受動的な方法です。サービスを受け取ることのメリットは作り手の苦労を知る必要がないことですが，サービスが打ち切られると手元に何も残らなくなるという問題があります。この講義の受講者は，第一，第二の方法でどうしてもうまくいかなかった場合に，私の方からサーバ環境によるサービスを提供しますので，最後の砦だと思っておいてください\footnote{ケチくさいこと言わずに全員に提供しろよー，と思うかもしれませんが，サーバが貧弱で全受講生が一斉にアクセスするとサーバが止まってしまう可能性が高いからです。ちなみにサーバを強くするためにはたくさんのお金がかかりますので，私個人のサービスでは限界があるのです。}。

% \subsection{導入方法1；ローカル環境での構築}
\subsection{Introduction Method 1; Building in a Local Environment}

% ローカル環境での構築は，R，RStudioの導入から入ります。ここは既にできている人もいるかもしれませんが，念のためバージョンが最新のものかどうかを確かめておいてください。
The construction in the local environment begins with the introduction of R and RStudio. Some people may have already done this, but just in case, please check if the version is the latest one.
% ローカル環境構築の方法がわからない人は，「RStudio」「インストール」などのキーワードでウェブ検索し，自分の環境にあった例を見つけてそれをみながら進めると良いでしょう。
For those who do not know how to set up a local environment, it is recommended to web search using keywords such as "RStudio" and "installation", and proceed while looking at examples suitable for your environment.
% 参考までにいくつかあげておきますと，新しいものでは，高知工科大学柳井先生のこちら\url{http://yukiyanai.github.io/jp/resources/}とか，Quiitaの記事\url{https://qiita.com/hujuu/items/ddd66ae8e6f3f989f2c0}が良いでしょう。
For your reference, some of the newer ones include Professor Yanai of Kochi University of Technology\\url{http://yukiyanai.github.io/jp/resources/} and articles on Quiita\\url{https://qiita.com/hujuu/items/ddd66ae8e6f3f989f2c0}.
% 書籍では，手前味噌で恐縮ですが\textcite{kosugi2019}などが良いでしょう\footnote{インストールに際しては，RStudio DesktopのFree editionを選んでください。RStudio CloudやRStudio Serverは別のものです。}。
In books, although it's self-praise, \textcite{kosugi2019} etc., might be good\footnote{When installing, please choose the Free edition of RStudio Desktop. RStudio Cloud and RStudio Server are different things.}.


% 次に確率的プログラミング言語，Stanの導入を行います。StanをRから使うためのパッケージとして\texttt{cmdstanr}と\texttt{rstan}という2種類があるという話はしました。かつては\texttt{rstan}のほうがよく使われていたのですが，最近は\texttt{cmdstanr}のほうが安定的に動作し，開発も盛んになっているので，2022年度後期からは\texttt{cmdstanr}をこの授業のデフォルトパッケージとしたいと思います。このパッケージの導入には，公式サイト\url{https://mc-stan.org/cmdstanr/articles/cmdstanr.html}を参考にしましょう。「cmdstanr インストール」で検索すると色々な紹介サイトが出てきますので，自分と同じ環境でなるべく日付の新しいものから参考にしていくといいでしょう。
Next, we will introduce the probabilistic programming language, Stan. We discussed that there are two types of packages to use Stan from R, \texttt{cmdstanr} and \texttt{rstan}. While \texttt{rstan} used to be more commonly used, recently \texttt{cmdstanr} has been functioning more stably and development is also becoming more active, so I would like to make \texttt{cmdstanr} the default package for this class from the second half of the fiscal year 2022. For the introduction of this package, please refer to the official site \url{https://mc-stan.org/cmdstanr/articles/cmdstanr.html}. If you search for "cmdstanr install", various introduction sites will come up, so it's a good idea to refer to the ones that are as new as possible in the same environment as yourself.

% インストールにあたっては，Stanがコンパイルされるときに利用するC++言語環境が必要です。これはMacOSであればXcodeの導入が，WindowsであればRtoolsでの導入が必要になります。
For the installation, a C++ language environment is required when Stan is compiled. If you are using MacOS, you will need to introduce Xcode, and if you are using Windows, you will need to introduce Rtools.
% \paragraph*{MacOSユーザの場合}AppStoreでXcodeと検索し，XcodeをインストールすればOKです。
\paragraph*{For MacOS Users} Simply search for Xcode in the AppStore and install it.

% \paragraph*{Windowsユーザの場合}RtoolsというR周辺のツールをインストールする必要があります。\url{https://cran.r-project.org/bin/windows/Rtools/}にいって，自分のRのバージョンにあったRToolsをインストールしてください。
\paragraph*{For Windows users} You need to install a tool called Rtools for R. Please go to \url{https://cran.r-project.org/bin/windows/Rtools/}, and install the Rtools version that matches your R version.
% この他にもWindowsユーザは導入にあたって色々障壁にあたることがありますので，こちらのサイト\url{https://norimune.net/3609}なども参考にしてみてください。
In addition to this, Windows users may encounter various obstacles when introducing it, so please refer to sites like this one: \url{https://norimune.net/3609}.
% 2022-09-18 17:35:05 rstanを第二候補に
% StanをRから使えるようにするものとして，\texttt{rstan}パッケージがあります。
% これもOSごとに導入の方法が違いますので，公式サイト\url{https://github.com/stan-dev/rstan/wiki/RStan-Getting-Started-(Japanese)}をみて導入にチャレンジしましょう。このサイトは「rstan getting started japanese」などで検索すると見つけることができます。

% \subsection{導入方法2；仮想環境での構築}
\subsection{Installation Method 2: Setting up in a Virtual Environment}

% 第二のルート，仮想環境を構築する場合ですが，これについては我らが国里先生が大変詳しいサイトを作ってくださっています。
The second option is to set up a virtual environment, and for this, our very own Professor Kuni has kindly created a very detailed site.

% 日本心理学会のチュートリアルワークショップ，「再現可能な日本語論文執筆入門：jpaRmdで実現する再現可能で低コストな日本語論文執筆のはじめの一歩」で使われた時の資料集がそれで，このサイト\url{https://ykunisato.github.io/jpa2021-tws-jpaRmd/}の事前準備編をよく読むと導入できます。通信環境にもよりますが，慣れているとものの10数秒でシステム導入できるという優れものです。面倒な設定が怖い人は，是非試してみてください。
The set of materials used in the tutorial workshop of the Japanese Psychological Association, "Introduction to Reproducible Japanese Manuscript Writing: The First Step to Realizing Reproducible and Low-Cost Japanese Manuscript Writing with jpaRmd," will introduce it well if you read the pre-preparation section of this site \url{https://ykunisato.github.io/jpa2021-tws-jpaRmd/}. Depending on your communication environment, it is an excellent system that can be introduced in just tens of seconds once you get used to it. If you are afraid of complicated settings, please try it.


% \section{導入方法3；外部サーバの利用}
\section{Introduction Method 3; Use of External Servers}

% 最近，Google社がオンラインで利用できる分析環境，Google Colaboratoryを提供してくれています。
Recently, Google has been providing an online analysis environment, Google Colaboratory, that you can use.
% これはGoogle社が提供する，ブラウザで実行できるプログラミング環境です。Googleが教育，研究用に提供しているもので，90分間は無料で利用できます。90分を超える場合でも，書いたプログラムを保存しておけば，新しいセッションを開始することで続けて利用できます。これを用いる利点は，どのユーザにも同じ環境を提供できることです。
This is a programming environment that can be run in a browser, provided by Google. This is something Google offers for education and research, which can be used for free for up to 90 minutes. Even if you exceed 90 minutes, you can continue to use it by saving the program you wrote and starting a new session. The advantage of using this is that it can provide the same environment to all users.

% 基本言語環境はPythonで提供されており，Jupyter Notebookを使いますが，次のアドレス(に含まれるコード)を使えば言語環境をRに変えて利用することができます。
The basic language environment is provided in Python and uses Jupyter Notebook, but you can change the language environment to R and use it by using the code contained in the following address.
\begin{center}
	\url{https://colab.research.google.com/notebook#create=true&language=r}
\end{center}


% ここでウィンドウにRのコードを入力し，実行していきます。\texttt{cmdstanr}や\texttt{rstan}もインストールできます。 
Input the R code into the window here and execute it. You can also install cmdstanr or rstan.
% インストールに際して，\verb|install.packages("rstan")|のように入力することもできますが，この方法ですとインストールに随分と時間がかかります。
When installing, you can also type \verb|install.packages("rstan")|, but this method takes quite a lot of time to install.
% そこで少し例外的ですが，システムコマンドを直接利用して，\verb|system("apt install -y r-cran-rstan")|のようにすることで，大幅にスピードを短縮することができます。
This is somewhat unconventional, but by directly using the system command, such as \verb|system("apt install -y r-cran-rstan")|, it's possible to significantly reduce the time required.

% \subsection{導入後の確認}
\subsection{Confirmation after Introduction}
% インストールが終わったらサンプルコードを実行して「動くかどうか」を確かめてみてください。サンプルコードはGetting started with CmdStanRのサイト，あるいはRStan Getting Startedのサイトに掲載されています。
Once the installation is complete, please run the sample code to check if it works or not. The sample code can be found on the 'Getting started with CmdStanR' website, or on the 'RStan Getting Started' website.
% どちらのパッケージから実行しても，文字列がズラズラっと出力されればとりあえず動いたと思っていただいて結構です。赤い文字やエラーなどが表示されうる，あるいは全く表示されない場合は，インストールがうまくいっていないことを疑いましょう。上に戻って，丁寧に説明に目を通して一歩ずつ進めてください。機械に命令することですので，自分勝手にショートカットしたり変更したりしないことが重要です。
Regardless of the package you run, as long as a string of text is output, you can assume it's working fine. If red letters or errors appear, or nothing is displayed at all, then you should suspect that the installation did not go well. Please go back, read the instructions carefully and proceed step by step. Since this is about giving commands to a machine, it's essential not to selfishly take shortcuts or make changes.
% \section{Stanを使ってみよう}
\section{Let's Try Using Stan}
% 具体的なStanコードの書き方や統計モデルへの応用については，本書の次の章から順に説明していきます。
The specific way to write Stan code and its application to statistical models will be explained in order from the next chapter of this book.
% それに先立って，全体的な注意をしておきたいと思います。具体的には，バージョンによる違い，パッケージによる違い，そして本書の以下の章で使う関数の紹介です。
Before we get started, I would like to make some general notices. Specifically, about differences due to versions, differences due to packages, and the introduction of functions that will be used in the following chapters of this book.

% \subsection{バージョンによる書き方の違い}
\subsection{Differences in Writing Depending on the Version}
% \verb|cmdstanr|と\verb|rstan|はパッケージの使い方以上に，それらが呼び出すStanのバージョンの違いがあります。上で解説したように，\verb|cmdstanr|のほうが最近は開発が進んでいて，より新しいStanの機能を使っていることになります。
There are differences in the versions of Stan called by `cmdstanr` and `rstan`, more so than the differences in how to use the packages. As explained above, `cmdstanr` has been more actively developed recently and therefore uses more recent features of Stan.
% バージョンが違っても，基本的な書き方などに違いはないのですが，バージョン2.26からStanの文法に変更が予定され，導入されました。
Even if the versions are different, there are no differences in basic writing methods, but from version 2.26, changes to Stan's grammar were planned and introduced.



% データXがN人分あったとします。数学的フォームでは，$X_1,X_2,X_3,\cdots, X_N$というようなものですね。これはコンピュータ上では\verb|X|と名付けられた配列，サイズ$N$というように考えます。これは，Stanでは次のように書くものでした。
Let's say we have data X for N people. In mathematical form, it would be $X_1,X_2,X_3,\cdots, X_N$. This is thought of as an array named \verb|X|, size $N$ on a computer. In Stan, we would write it like this.
\begin{lstlisting}[language=C++,label={code:PRE_01},caption={Stan2.26以前の書き方}]
int n[5];
real a[3, 4];
real<lower=0> z[5, 4, 2];
vector[7] mu[3];
\end{lstlisting}


% これらは上から順に，サイズ$5$の整数変数\verb|n|,サイズ$3 \times 4$の実数変数\verb|a|，下限$0$のサイズ$5\to,es4\times2$の配列\verb|z|，長さ$7$のベクトル\verb|mu|が3つ，ということを意味しています\footnote{ベクトルは数字をセット，人まとめとして演算するものです。最後の変数は，7つの数字のセットが3つという意味で，7つセットで1つのまとまりですよ，ということを明示的に宣言していることになります。}。
This, in order from above, means that there is an integer variable \verb|n| of size $5$, a real variable \verb|a| of size $3 \times 4$, an array \verb|z| of size $4 \times 2$ with a lower limit of $0$, and three vectors \verb|mu| of length $7$ \footnote{Vectors are used to operate on a set or a collection of numbers. The last variable implies that there are three sets of seven numbers, explicitly stating that each set of seven constitutes a single unit.}.

% この書き方が，Stanのバージョン3.32以降は次のような書き方になります。
This notation will change to the following style in versions 3.32 and later of Stan.
\begin{lstlisting}[language=C++,label={code:PRE_02},caption={Stan2.32以降の書き方}]
array[5] int n;
array[3, 4] real a;
array[5, 4, 2] real<lower=0> z;
array[3] vector[7] mu;
\end{lstlisting}

% このコード\ref{code:PRE_02}はさきほどのコード\ref{code:PRE_01}と同じ意味で，書き方が変わっただけです。一瞥してわかるように，サイズを先に\verb|array|という言葉で宣言しましょう，というだけですね。
This code \ref{code: PRE_02} has the same meaning as the previous code \ref{code: PRE_01}, only the way of writing has changed. As you can see at a glance, it's just about declaring the size first with the word \verb|array|.

% さて，今現在はStan2.26と3.32の間，過渡期になります。ですから，どちらのコードで書いてもエラーにはなりません。ただし，(ここからが少し面倒なのですが)
Well, currently, we are in a transition period between Stan 2.26 and 3.32. Therefore, writing in either code won't cause an error. However, (this is where it gets a bit troublesome)

\begin{itemize}
  \item RStudioのコードチェック機能は新しいコードの書き方に対応しておらず\footnote{RStudioのバージョン2022.07.2 Build 576で確認しました。}，新しい書き方をするとエディタ上でエラーの警告(赤いバッテンがつきます)が出ますし，チェックボタンを押してもSYNTAX ERRORが出ます。
  \item \verb|rstan|パッケージは使っているStanが古いこともあって，新しい書き方のコードをコンパイルしようとするとエラーになります。RStudioと同じ，SYNTAX ERRORが出ます。
  \item \verb|cmdstanr|パッケージは逆に，新しいStanの書き方を推奨していますので，「その書き方は古いよ」という警告を出してきます。Declaration of arrays by placing brackets after a variable name is deprecated and will be removed in Stan 2.32.0. Instead use the array keyword before the type. This can be changed automatically using the auto-format flag to stanc.という警告がでますが，これは「配列のカッコを変数の後ろに置く書き方はもうダメで，Stan2.32以降は無くなります。型の前にarrayと書くようにしてください。これはauto-formatフラグを立てることで自動的に変更するようになります。」という意味です。
\end{itemize}


% つまり，本書は\verb|cmdstanr|を推奨していますが，そうするとRStudioのエディタ上ではエラーが出て文法チェックができず，コンパイルした後でいざサンプリング，というときにエラーが出たりします。チェックをするには，コンパイルしたオブジェクトを使ってサンプリングをする前に，\verb|model$check_syntax()|関数を実行する必要があります。
In other words, while this book recommends using \verb|cmdstanr|, doing so can lead to errors in the RStudio editor and make it impossible to check the syntax, resulting in errors when you try to sample after compiling. To check, you need to run the \verb|model$check_syntax()| function before sampling with the compiled object.

% \verb|rstan|パッケージを使うとスマートに文法チェックもできていいのですが，\verb|rstan|そのものが不安定ですし，今後徐々に開発・発展をしなくなっていくことが明らかですので，どこかで新しい方に舵を切る必要があります。
Using the \verb|rstan| package lets you smartly do grammar checks, but \verb|rstan| itself is unstable, and it's clear that it will gradually stop developing and evolving, so you'll need to steer in a new direction at some point.

% 本書のコードは新しい書き方に統一していますが，最初のうちは，そして今しばらくは，実際にコードを書くときは古い方で書いたほうが便利かもしれません。
The code in this book has been unified into a new style, but initially, and for a while, it might be more convenient to write using the old style when actually writing code.
% \subsection{パッケージによる指示と出力の違い}
\subsection{Differences in Instructions and Outputs by the Package}
% さて今度はパッケージごとの違いを見ていきましょう。
Now, let's take a look at the differences between each package.
% \subsubsection{指示の仕方の違い}
\subsubsection{Differences in Giving Instructions}
% まずは実行方法です。Stanでは事後分布からの乱数を生成しますが，それにあたって「乱数の数」「ウォームアップ期間の長さ」「チェインの数」などオプショナルに指定できるものがあります。これは見てもらったほうが早いかと思いますので，同じことをそれぞれのパッケージで実行するときにどのように指定方法を変えるかを見てみましょう。まずは\verb|rstan|パッケージの書き方からです。
First, let's talk about execution methods. In Stan, we generate random numbers from posterior distributions. There are things that we can specify optionally, such as the "number of random numbers," the "length of the warm-up period," and the "number of chains." It's probably faster to see for yourself, so let's look at how to change the specification when executing the same thing in each package. Let's start with the \verb|rstan| package.

\begin{lstlisting}[language=C++,label={code:PRE_03rstan},caption={rstanのスタイル}]
fit <- rstan::sampling(model,
  data = dataSet,
  chains = 4,
  iter = 6000,
  warmup = 1000
)
\end{lstlisting}

% ここで指定しているのは，データを\verb|dataSet|として設定し，同時に4本のチェインを発生させ，6000個の乱数を作って，そのうち1000個はまだ機械が温まっていないので候補から除外する，というものです。複数のチェインを作ること，温まっていない部分を捨てる理由などは後ほど説明しますが，結果的に合計20000個の乱数が作られることを知っておいてください。この数字は$(6000-1000)\times4 = 20000$という計算からでてきます。また，乱数の発生は並列計算させた方が良いので，\verb|rstan|パッケージを使う場合は事前に，\verb|options(mc.cores = parallel::detectCores())|という一行を入れておきます。
What is being specified here is to set the data as "dataSet", generate four chains at the same time, create 6000 random numbers, and exclude 1000 of them from consideration because the machine is still warming up. The reasons for creating multiple chains and discarding the not yet warmed up part will be explained later, but please be aware that a total of 20,000 random numbers are generated in the end. This number comes from the calculation of $(6000-1000)\times4 = 20000$. Also, because it is better to perform the generation of random numbers in parallel calculations, if you are using the rstan package, you should insert a line of "options(mc.cores = parallel::detectCores())" in advance.

% さて，同じことを\verb|cmdstanr|パッケージでやると次のようになります。
Now, if you do the same thing with the \verb|cmdstanr| package, it will look like this.
\begin{lstlisting}[language=C++,label={code:PRE_04cmd},caption={cmdstanrのスタイル}]
fit <- model$sample(
  data = dataSet,
  chains = 4,
  parallel_chains = 4
  iter_sampling = 5000,
  iter_warmup = 1000,
)
\end{lstlisting}

% ここにあるように，並列するチェインの数をオプションに書き込む必要があります。またサンプルの数は捨てる部分を別にして(含めずに)計算させることができます。このコードで20000個のサンプルが得られます。細かいことですが，多少の違いがあるわけです。
As shown here, you need to write the number of parallel chains in the options. Also, the number of samples can be calculated without including the part to be discarded. This code will give you 20,000 samples. It's a minor thing, but there is a bit of a difference.

% \subsubsection{出力の違い}
\subsubsection{Difference in Output}
% さて，今度は出力結果の使い方についての注意です。\verb|rstan|パッケージは，出力結果を\verb|stanfit|オブジェクト型という形にします。\verb|rstan|のもっているさまざまな関数，たとえば要約した結果の出力やプロットなどは，\verb|stanfit|オブジェクトだとわかるとそれに対応した出力に変えて表現してくれるわけです。\verb|cmdstanr|の出力はまた別物\footnote{\verb|class|関数で\verb|rstan|の出力がどういうクラスなのかを表示させると，\verb|stanfit|と出てくるのですが，\verb|cmdstanr|の出力を同様にチェックすると\verb|"CmdStanMCMC" "CmdStanFit"  "R6"|という答えが返ってきます。クラスというのはデータの形を定義する方法で，オブジェクト指向プログラミングに必要な概念なのですが，ここではややこしすぎるのでパス。違うということだけわかってもらえれば結構です。RStudioのばあいはEnvironmentタブを見るだけでも違いがわかると思います。\verb|rstan|パッケージの出力はちゃんとしたクラス(Formla Class stanmodel)なのですが，\verb|cmdsatnar|パッケージの出力はEnvironment(グローバル環境の要素)として剥き出しの出力が出ている感じです。}です。
Now, let's consider some points about how to use the output results. The rstan package formats its output results into something called a stanfit object type. Various functions in rstan, such as the ones outputting summarized results or plots, will convert and display them in a corresponding manner once they recognize it as a stanfit object. The output from cmdstanr is a separate matter. If you use the class function to display what class the output from rstan is, it will return as stanfit. If you check the output from cmdstanr in the same way, you'll get "CmdStanMCMC", "CmdStanFit", "R6". A class is a method to define the form of data and a concept needed for object-oriented programming, but it's too complicated to explain here. All you need to know is the two are different. I believe you can tell the difference just by looking at the Environment tab in RStudio. While the output from the rstan package is a properly classified data(Formula Class stanmodel), the output from the cmdstanr package appears raw, as an element in the global environment.

% 同じモデルをそれぞれのパッケージで出力させて，違いを見てみましょう。まずは\verb|rstan|パッケージの出力から。
Let's generate the same model using each package and see the differences. First, let's start with the output from the \verb|rstan| package.

\begin{SCrstan}{rstanパッケージの出力}{outRstan}
	\begin{verbatim}
  Inference for Stan model: ttest01.
  4 chains, each with iter=6000; warmup=1000; thin=1; 
  post-warmup draws per chain=5000, total post-warmup draws=20000.
  
         mean se_mean   sd   2.5%    25%    50%    75%  97.5% n_eff Rhat
  mu1   59.96    0.04 5.58  49.01  56.33  59.99  63.60  70.91 16043    1
  mu2   40.04    0.05 5.62  28.98  36.38  40.03  43.67  51.17 14330    1
  sig   15.54    0.03 3.12  10.86  13.35  15.08  17.23  22.92 12118    1
  lp__ -51.58    0.02 1.36 -55.12 -52.18 -51.24 -50.60 -50.05  7115    1
  
  Samples were drawn using NUTS(diag_e) at Tue Nov 15 09:50:17 2022.
  For each parameter, n_eff  a crude measure of effective sample size,
  and Rhat is the potential scale reduction factor on split chains (at 
  convergence, Rhat=1).
\end{verbatim}
\end{SCrstan}


% \verb|rstan|パッケージが出力しているのは次のような情報です。
The information that the \verb|rstan| package is outputting is as follows.
\begin{itemize}
  \item モデル名，チェイン数，反復回数，最終的に得られたMCMCサンプル数(draws)などの説明
  \item パラメータの事後分布の記述統計量，すなわちベイズ推定による事後分布の情報。順に平均値(\verb|mean|)，平均値の標準誤差(\verb|se_mean|)，標準偏差(\verb|sd|)，2.5\%,25\%,50\%,75\%,97.5\%パーセンタイル。
  \item MCMCサンプルについての特徴量。有効サンプルサイズ(\verb|n_eff|)とRhat(\verb|Rhat|)
  \item サンプリングに関するあとがき
\end{itemize}

% それぞれの内容についてはあとの章で説明するとして，続けて\verb|cmdstanr|パッケージの出力をみてみます。
We will explain each topic in the following chapters, but for now, let's take a look at the output of the cmdstanr package.
\begin{SCcmd}{cmdstanrパッケージの出力}{outCMDstan}
\begin{verbatim}
  variable   mean median   sd  mad     q5    q95 rhat ess_bulk ess_tail
  lp__ -51.56 -51.20 1.34 1.07 -54.18 -50.12 1.00     7799    10071
  mu1   59.99  60.00 5.63 5.47  50.82  69.16 1.00    16325    12764
  mu2   40.05  40.05 5.48 5.18  31.09  49.05 1.00    16511    12466
  sig   15.48  15.03 3.07 2.75  11.41  21.07 1.00    14489    12137
\end{verbatim}
\end{SCcmd}


% 同じような情報ですが，ちょっと違うところもあるようです。前から順に，事後分布の平均値(\verb|mean|)，中央値(\verb|median|)，標準偏差(\verb|sd|)，平均絶対偏差\footnote{\keyterm{中央値絶対偏差}{中央値ぜったいへんさ}{Median Absolute Deviation}とは，中央力のばらつきを表す指標で，各データ点の中央値からの偏差の絶対値をとり，その中央値を計算したものです。標準偏差とは違う分布の幅を表現する方法で，中央値を使って計算しますから外れ値に強いという特徴があります。}(\verb|mad|)，5\%,95\%パーセンタイルなど分布の特徴\footnote{\verb|rstan|パッケージは2.5\%から97.5\%のパーセンタイル，すなわち囲まれる区間の面積が95\%です。\verb|cmdstanr|パッケージは5\%と95\%のパーセンタイル，すなわち囲まれる区間の面積が90\%です。なぜだかわかりませんが，ちょっとした違いがありますね。}と，Rhatや有効サンプルサイズに関する情報(\verb|ess_bulk,ess_tail|)が示されています。しかし，まえがき・あとがきのようなものがなく，結果だけがシンプルに出力されていますね。
The information seems similar, but there are slightly different parts too. In order from the front, the mean, median, standard deviation, median absolute deviation, and characteristics such as the 5%, 95% percentiles of the posterior distribution, as well as information about Rhat and effective sample size (ess_bulk, ess_tail) are shown. The median absolute deviation is an indicator of spread around the median, it is calculated by taking the absolute deviation from the median of each data point and calculating its median. This is a different way to express the width of a distribution than the standard deviation and has the characteristic of being robust to outliers because it is calculated using the median. However, there are no prologues or epilogues, and only the results are simply output. As for the package differences, the rstan package shows percentiles from 2.5% to 97.5%, that is, the area of the enclosed interval is 95%. The cmdstanr package shows 5% and 95% percentiles, that is, the area of the enclosed interval is 90%. I don't know why, but there are small differences.

% 出力としては\verb|rstan|パッケージの方が丁寧で良いかもしれません。さて，\verb|cmdstanr|の出力は実はcsvファイルとして一時フォルダに保存されているので，これを\verb|rstan|パッケージの関数を使って\verb|stanfit|オブジェクトに変換することができます。関数の使い方の例は次のようになります。
The \verb|rstan| package may be better for the output. Now, the output of \verb|cmdstanr| is actually saved as a csv file in a temporary folder, so it can be converted to a \verb|stanfit| object using a function from the \verb|rstan| package. An example of how to use the function is as follows.

\begin{lstlisting}[language=C++,label={code:PRE_04},caption={cmdstanrの出力をrstanのオブジェクトに変える}]
fitR2 <- fitC$output_files() |> rstan::read_stan_csv() 
\end{lstlisting}


% コードの中身ですが，まず\verb|fitC|と書いてあるのが\verb|cmdstanr|の出力オブジェクトです。このオブジェクトがさまざまな情報を持っているという形になっており，ここから出力ファイルの場所を聞き出すのが\verb|output_files()|です。このファイルを\verb|rstan|パッケージの\verb|read_stan_csv()|関数に渡してやると，csvファイルを読み込んでrstanのオブジェクトに変えてくれます\footnote{当然のことながら，\verb|rstan|パッケージがないとこの関数も呼び出せませんので，こちらもインストールしておく必要があります。なおこのコードの，\texttt{|>}は\verb|tidyverse|パッケージのパイプ演算子，\verb|%>%|と同じ機能を持つRの演算子で，ネイティブパイプと呼ばれています。ネイティブパイプはR4.1.0以降に導入されたものです。}。変えたものをここでは，\verb|fitR2|オブジェクトに保存しています。この方法を使うと，先ほどの\verb|cmdstanr|の出力が\verb|rstan|の出力と同じになります(同じなので再掲しません。)。
Regarding the content of the code, first of all, \verb|fitC| written here is the output object of \verb|cmdstanr|. This object carries various information and \verb|output_files()| is used to extract the location of the output file from it. If you pass this file to the \verb|read_stan_csv()| function of the \verb|rstan| package, it will read the csv file and convert it into an rstan object\footnote{Of course, if there is no \verb|rstan| package, you cannot call this function, so you need to install it as well. This code's \texttt{|>} is an R operator with the same function as the pipe operator, \verb|%>%| of the \verb|tidyverse| package, and it's called a native pipe. The native pipe was introduced from R4.1.0 onwards.}. The converted object is saved here as \verb|fitR2|. By using this method, the output of \verb|cmdstanr| becomes the same as the output of \verb|rstan| (the same as there is no re-posting).

% このコードを覚えておけば，\verb|rstan|パッケージ向けに書かれたものであっても，\verb|cmdstanr|パッケージで実行したあと同じ関数を適用できるので，使いやすいですね。
Remember this code, even if it's written for the rstan package, you can apply the same functions after running it in the cmdstanr package, which is convenient.

% \subsection{本書で利用する準備関数}
\subsection{Preparation Functions Used in this Book}
\label{preparedFunctions}
% どちらのパッケージであっても，欲しいものは事後分布からの乱数をデータセットにしたものであり，そのデータセットさえ持っていれば，既存の関数を使わなくとも自分で工夫して加工すれば良いでしょう。
Regardless of which package you use, what you want is a dataset created from random numbers from the posterior distribution. If you have that dataset, you do not need to use existing functions and can devise and process it yourself.

% そこで本書では，MCMCサンプルを抜きだしてデータセットに変換する関数を作り，これを利用して話を進めることにします。
So in this book, we will create a function to extract the MCMC sample and convert it into a dataset, and proceed with the discussion using this.
% ここではその関数の中身を解説しておきます。
Here, I will explain the contents of that function.
% \subsubsection{MCMCサンプルをデータフレームにする関数}
\subsubsection{Function to Convert MCMC Samples to a Data Frame}
% まずはMCMCサンプルをデータセットにして渡してくれる関数です。この関数は伴走サイトのコードの冒頭に必ず含まれており，以後はこれを使って結果の要約を示すといった使い方をしていますので，中身を知っておいてもらった方が良いかもしれません。自分の使っているパッケージの方だけでも目を通してください。この関数を経由すると，出てくるオブジェクトは同じ形式になっているところがミソです。
Firstly, this is a function that takes MCMC samples as a dataset. This function is always included at the beginning of the co-running site's code, and afterwards, it is used to show a summary of the results, so it might be good to know what's inside it. Please review at least the package you are using. The key point is that the object that comes out via this function has the same format.

% \paragraph{stanfitオブジェクト(rstanパッケージ)の場合}
\paragraph{In the case of the stanfit object (rstan package)}
% \verb|rstan|パッケージを使って\verb|stanfit|オブジェクトとしてMCMCサンプルを得た場合，そのオブジェクトから乱数部分だけを抜き出す関数は\verb|rstan::extract|関数です。これで取り出したものをデータフレーム型にして返す関数として，次のようなものを作りました。
When using the `rstan` package to obtain MCMC samples as a `stanfit` object, the function to extract only the random number part from that object is the `rstan::extract` function. I created the following function to extract this and return it as a dataframe type.

\begin{lstlisting}[language=C++,label={code:PRE_05MCMC2dfRSTAN},caption={MCMCtoDF関数(rstan版)}]
MCMCtoDF <- function(fit) {
  fit %>%
    rstan::extract() %>%
    as.data.frame() %>%
    tibble::as_tibble() %>%
    tibble::rowid_to_column("iter") %>%
    dplyr::select(-lp__) %>%
    tidyr::pivot_longer(-iter) -> MCMCsample
  return(MCMCsample)
}
\end{lstlisting}

% \paragraph{コード解説}
\paragraph{Code Explanation}
\begin{description}
	\item[1行目] 関数を作る宣言。引数として\verb|stanfit|オブジェクトをとります。
	\item[2行目] 引数で引き受けたオブジェクトを加工していきます。
	\item[3行目] まずはMCMCサンプルを抜き出します。
	\item[4行目] data.frame型にします。
	\item[5行目] tibble型にします。別にしなくてもいいんですが，著者の好みです。
	\item[6行目] 行番号を表す変数\verb|iter|を作ります。
	\item[7行目] 変数の中から\verb|lp__|を除外します。
	\item[8行目] 行番号変数は残して，tidyなデータにし，\verb|MCMCsample|というオブジェクトに保存します
	\item[9行目] 戻り値として\verb|MCMCsample|を返します。
\end{description}



% \paragraph{cmdstanrパッケージの場合}
In the case of the cmdstanr package:
% \verb|cmdstanr|パッケージを使ってMCMCサンプルを得た場合，そのオブジェクトから乱数部分だけを抜き出す関数は\verb|draws|関数であり，これをオブジェクトに直接作用させます。これで取り出したものをデータフレーム型にして返す関数として，次のようなものを作りました。
The function to extract only the random part from the object obtained from MCMC samples using the cmdstanr package is the draws function, and it applies directly to the object. I've created the following function to return what's extracted in dataframe format.


\begin{lstlisting}[language=C++,label={code:PRE_05MCMC2dfCMDstan},caption={MCMCtoDF関数(cmdstanr版)}]
MCMCtoDF <- function(fit) {
  fit$draws() %>%
    posterior::as_draws_df() %>%
    tibble::as_tibble() %>%
    dplyr::select(-lp__, -.draw, -.chain, -.iteration) %>%
    tibble::rowid_to_column("iter") %>%
    tidyr::pivot_longer(-iter) -> MCMCsample
  return(MCMCsample)
}
\end{lstlisting}

% \paragraph{コード解説}
\paragraph{Code Explanation}
\begin{description}
	\item[1行目] 関数を作る宣言。引数として\verb|stanfit|オブジェクトをとります。
	\item[2行目] 引数で引き受けたオブジェクトに\verb|draws|関数を作用させ，サンプルだけ取り出します。
	\item[3行目] 取り出したサンプルをデータフレーム型にする\verb|posterior|パッケージの\verb|as_draws_df|関数を適用します。
	\item[4行目] tibble型にします。別にしなくてもいいんですが，著者の好みです。
	\item[5行目] 変数の中から\verb|lp__|や他の隠し変数を除外します。
	\item[6行目] 行番号を表す変数\verb|iter|を作ります。
	\item[7行目] 行番号変数は残して，tidyなデータにし，\verb|MCMCsample|というオブジェクトに保存します
	\item[8行目] 戻り値として\verb|MCMCsample|を返します。
\end{description}



% \paragraph{MCMCtoDF関数の結果}
Sorry, but the text you provided is a piece of script code or a method name (MCMCtoDF function results). For translations, please provide the actual Japanese text.
% 出力結果は，どちらも同じく以下のような形式になります。
The output results will be in the same format for both.
\begin{Rscreen}{MCMCサンプルをデータセットにしたもの}{outPREP_01}
  \begin{verbatim}
# A tibble: 60,000 × 3
    iter name  value
   <int> <chr> <dbl>
 1     1 mu1    60.3
 2     1 mu2    41.0
 3     1 sig    12.5
 4     2 mu1    62.9
 5     2 mu2    37.6
 6     2 sig    11.5
 7     3 mu1    68.7
 8     3 mu2    46.4
 9     3 sig    13.7
10     4 mu1    67.7
# … with 59,990 more rows
# ℹ Use `print(n = ...)` to see more rows
  \end{verbatim}
\end{Rscreen}

% ここで\verb|iter|とあるのはMCMCのステップ番号，\verb|name|とあるのが変数名，\verb|value|とあるのがその値です。tidyな形になっていますから，このまま分析や描画に用いることができます。
What the \verb|iter| here refers to is the step number of MCMC, what's referred to as \verb|name| is the variable name, and what's referred to as \verb|value| is its value. Since it's in a tidy form, you can use it as is for analysis and drawing.

% \subsubsection{MCMCサンプルの結果を要約する関数}
\subsubsection{Function to Summarize the Results of MCMC Samples}
% MCMCの結果を分析するにあたっては，\verb|bayestestR|パッケージなど既存のパッケージを使うと便利でしょう。
When analyzing the results of MCMC, it may be convenient to use existing packages such as the \verb|bayestestR| package.
% ですが，こうしたパッケージはどんどん開発が進んでアップデートされたり，\verb|rstan|か\verb|cmdstanr|かで書き方が変わったりするなど，テキストで紹介するには難しいところがあります。そもそもMCMCサンプルを加工して記述統計を示したり，描画したりするわけですから，上で紹介したようにMCMCサンプルを抜き出すことができれば自力でもできるはず。ということで，自作の関数を用意しました。本書では以下，この関数をつかって解説しますので中身を知っておくと良いでしょう。
However, such packages are constantly being developed and updated, and the codes can change depending on whether they are written in 'rstan' or 'cmdstanr', making it difficult to introduce in a text. After all, we're processing MCMC samples to show descriptive statistics or make drawings, so as long as you can extract MCMC samples as introduced above, theoretically one should be able to do it on their own. That said, I have prepared a custom function. In this book, I will explain using this function, so it would be good to be familiar with what's inside.

% \paragraph{MAP推定関数}
\paragraph{MAP Estimation Function}
% 確率分布の特徴を報告するときに，その確率分布に従う乱数を使って，期待値や中央値を計算するのは簡単です。記述統計としての平均値やパーセンタイルでよいからです。しかし確率分布の確率密度が最も高くなるところ，すなわち\keyterm{MAP推定値}{MAPすいていち}{MAP Estimaiton} の計算はちょっと難しいですね。というのも，関数であれば微分して極値を求めれば良いのですが，MCMCサンプルは関数ではなくて値なので，ヒストグラムを書くしかなく，最大の値を求める計算がむずかしいからです。
When reporting the characteristics of a probability distribution, it is easy to calculate expectations and medians using random numbers that follow that probability distribution. This is because descriptive statistics such as mean values and percentiles are sufficient. However, it is a bit difficult to calculate the most probable value of a probability distribution, i.e., the MAP (Maximum A Posteriori) estimation. This is because if it is a function, you can differentiate it and find the extreme value, but MCMC samples are not functions but values, so all you can do is draw a histogram, and it is difficult to calculate the maximum value.

% でも大丈夫，Rの描画関数を駆使して対応することができます。MAP推定値を計算する関数は，次のように書きます\footnote{このコードは関西学院大学社会学部の清水裕士先生に教えてもらったものです。記して謝意を表します。}
But it's okay, we can deal with it by making full use of R's drawing functions. The function to calculate the MAP estimate is written as follows \footnote{This code was taught to me by Mr. Hiroshi Shimizu of the Kwansei Gakuin University Faculty of Sociology. I would like to express my gratitude.}

\begin{lstlisting}[language=C++,label={code:PRE_06MAP},caption={MAP推定する関数のコード}]
map_estimation <- function(z) {
  density(z)$x[which.max(density(z)$y)]
}
\end{lstlisting}

% これはRの\verb|density|関数でヒストグラムに密度関数を当てがい(カーネル密度推定)，その関数の特徴から値を返すようになっています。
"This is programmed to apply a density function (Kernel Density Estimation) to a histogram using R's \verb|density| function, and return values from the characteristics of that function."

% \paragraph{MCMCサンプルの要約関数}
\paragraph{Summary Function of MCMC Samples}
% 先ほどのMAP推定する関数を含めつつ，またMCMCサンプルをデータフレームにする関数を使いつつ，MCMCサンプルの要約を表示する関数を次のように作りました。
I created a function to display a summary of the MCMC (Markov chain Monte Carlo) samples, while using the function to convert MCMC samples into a data frame, and also including the function to estimate MAP (Maximum A Posteriori) that we made earlier.
\begin{lstlisting}[language=C++,label={code:PRE_07summary},caption={MCMCの要約を報告するコード}]
MCMCsummary <- function(MCMCsample) {
  MCMCsample %>%
    dplyr::group_by(name) %>%
    dplyr::summarise(
      EAP = mean(value),
      MED = median(value),
      MAP = map_estimation(value),
      SD = sd(value),
      U95 = quantile(value, prob = 0.975),
      L95 = quantile(value, prob = 0.025)
    ) %>%
    mutate(across(where(is.numeric), ~ num(., digits = 3)))
}
\end{lstlisting}

% \paragraph{コード解説}
\paragraph{Code Explanation}
\begin{description}
	\item[1行目] 関数を作る宣言。引数として\verb|MCMCtoDF|関数の戻り値をとります。
	\item[2行目] 引数で引き受けたオブジェクトを加工していきます。
	\item[3行目] 変数名でデータをグループ化します。
	\item[4-10行目] 記述統計の計算を通じて，EAP,MED,MAP,SD,95\%区間を返します。
	\item[11]行目 出力する数値データは，少数下3桁まで表示させるようにします。
\end{description}


% この関数は，先ほどの\verb|MCMCtoDF|とセットで使い,次のような結果を得ます。
This function is used in conjunction with the previous \verb|MCMCtoDF|, and yields the following results.
\begin{Rscreen}{MCMCsummaryの結果}{outPREP_02}
	\begin{verbatim}
> fit %>% MCMCtoDF() %>% MCMCsummary()
# A tibble: 3 × 7
  name        EAP       MED       MAP        SD       L95       U95
  <chr> <num:.3!> <num:.3!> <num:.3!> <num:.3!> <num:.3!> <num:.3!>
1 mu1      59.897    59.888    58.867     5.547    48.672    70.847
2 mu2      39.958    39.935    39.748     5.523    28.973    50.943
3 sig      15.519    15.062    14.646     3.074    10.877    22.788
  \end{verbatim}
\end{Rscreen}

% ここにあるように，各変数のEAP推定値，MED推定値，MAP推定値，95\%確信区間が表示されます\footnote{この区間はコードから明らかなように，\verb|quantile|関数で計算しています。下から2.5\%，上から2.5\%を取り除くと95\%の区間が残ることになりますが，このように分布の両端を均等に切った区間を厳密には\keyterm{等裾区間}{とうすそくかん}{Equal-tailed interval;ETI}というのでした。これに対して，分布の最も信頼できる部分，かつ分布の大部分をカバーし，区間内部のすべての点が，区間外の任意の点よりも高い密度を持っているように推定するものを\keyterm{最高密度区間}{さいこうみつどくかん}{Highest-Density Interval;HDI}と呼びます。これらは左右対称の単峰分布であれば同じになるのですが，事後分布の形がぐちゃっとしている時は必ずしもそうなりません。両者の違いについては\textcite{Maeda2017}を参考にしてください。また，HDIを出力する関数は\verb|bayestestR::hdi()|です。}。
As shown here, the EAP, MED, and MAP estimates and 95\% credible intervals for each variable are displayed\footnote{As obvious from the code, these intervals are calculated using the \verb|quantile| function. If you remove 2.5% from the bottom and 2.5% from the top, you are left with a 95\% interval, which is strictly called an Equal-tailed interval (ETI). In contrast, the Highest-Density Interval (HDI) is called what is estimated to be the most reliable part of the distribution, covering most of the distribution, and all points inside the interval have a higher density than any point outside the interval. These are the same if they are a symmetric unimodal distribution, but not necessarily so if the shape of the posterior distribution is messy. For the differences between the two, please refer to \textcite{Maeda2017}. Also, the function to output the HDI is \verb|bayestestR::hdi()|.}.

% これらの推定値や，そもそも事後分布が何をどのようなことを表しているのか，といった点については今後の授業で説明していきます。以下は結果をしれっとこの関数で表現していきますので，何をやっているのかわからなくなったらここに立ち戻ってきてください。
We will explain these estimates and what the posterior distribution represents in future classes. The results are casually expressed with this function below, so if you lose track of what is being done, please come back here.
\end{document}
